\documentclass[12pt,letterpaper]{article}

% ---- Packages ----
\usepackage[utf8]{inputenc}
\usepackage[T1]{fontenc}
\usepackage{amsmath,amssymb}
\usepackage{newtxtext,newtxmath}  % Times New Roman text + math (replaces mathptmx)
\usepackage[margin=1in]{geometry}
\usepackage{setspace}
\usepackage{titlesec}
\usepackage{enumitem}
\usepackage{hyperref}
\usepackage{microtype}
\usepackage{booktabs}

% ---- Section formatting ----
\titleformat{\section}{\large\bfseries\sffamily}{\thesection.}{0.5em}{}
\titleformat{\subsection}{\normalsize\bfseries\sffamily}{\thesubsection}{0.5em}{}
\titleformat{\subsubsection}{\normalsize\bfseries}{\thesubsubsection}{0.5em}{}

% ---- Spacing ----
\setstretch{1.15}
\setlength{\parskip}{0.5em}
\setlength{\parindent}{0em}

% ---- Custom commands ----
\newcommand{\rhoa}{\rho_{\!A}}
\newcommand{\Ka}{K_{\!A}}
\newcommand{\etaK}{\eta_{\!K}}

\begin{document}

\begin{center}
{\LARGE\bfseries Temperature-Dependent Bulk Modulus\\
as the Origin of Dark Energy and Dark Matter}\\[1em]
{\large Erik Otto}\\[0.5em]
{\normalsize\itshape Working Paper in the Dynamic Aether Mechanics Framework}
\end{center}

\vspace{1em}

\begin{abstract}
Paper~6 identifies dark energy as the thermal pressure of hot void Aether
and dark matter as the two-fluid first-sound correction from an outward
temperature gradient.\ This working paper explores an alternative thermal
landscape in which matter-containing regions are \emph{warmer} than
voids---heated by viscous dissipation of gravitational inflow---while voids
cool adiabatically through expansion.\ If the Aether bulk modulus $\Ka$
increases with decreasing temperature (as the first-sound speed does in
superfluid helium-4), the wave speed $c = \sqrt{\Ka/\rhoa}$ is higher in
cold voids than near warm mass.\ On galactic scales, the outward-increasing
$c$ produces dark matter through the acoustic metric ($\Delta\Ka/\Ka \sim
4 \times 10^{-6}$).\ On cosmic scales, the growing $\Ka$ contrast from a
cooling cascade drives accelerating expansion ($\Delta\Ka/\Ka \sim 2 \times
10^{-4}$).\ A single parameter $\etaK \sim 10^{-4}$~K$^{-1}$---matching
the measured temperature sensitivity of helium-4 at low
temperature---governs both phenomena.\ The cascade naturally produces the
phantom-to-quintessence crossing observed by DESI.\ The local Aether
temperature at the Sun's position ($\sim$2.75~K) matches
$T_{\mathrm{CMB}}$ to within 1\%.
\end{abstract}

\vspace{1em}

\bigskip
\noindent\textit{This working paper explores an alternative temperature model
within the Dynamic Aether Mechanics framework.\ It references results from
Paper~1 (General Overview), Paper~3 (Gravity), Paper~6 (Cosmology), and
Paper~6(a) (Cooling Feedback Model) of the main series.\ The ideas presented
here are under development and may be incorporated into the main papers in
revised form.}
\bigskip


% ============================================================
\section{Introduction}
% ============================================================

Paper~6 establishes that the Aether possesses three independently varying
properties---density, pressure, and temperature---decoupled by the mechanism
of permanent binding \cite{otto-overview}.\ In that framework, matter-containing
regions are \emph{cold} (thermal energy locked into mass by binding) and
cosmic voids are \emph{hot} (heated by stellar radiation and bulk-viscous
dissipation of the Hubble flow).\ This temperature landscape produces dark
energy (thermal pressure in hot voids driving expansion) and dark matter
(the two-fluid first-sound correction from the outward temperature gradient)
\cite{otto-cosmology}.

This paper asks: \emph{what if the temperature gradient is reversed?}

Three observations motivate the question:

\begin{enumerate}[leftmargin=2em]
  \item \textbf{Matter creation was early.}\ Permanent binding occurred
    primarily in the first few billion years of cosmic history.\ The thermal
    deficit it created was enormous but finite, and stellar radiation plus
    viscous dissipation have been returning energy to the near-mass Aether
    for over ten billion years.\ The initial cooling from binding may be
    long since compensated.

  \item \textbf{Gravitational inflow creates large velocity gradients.}\
    Near a Milky-Way-sized galaxy, the inflow velocity gradient
    $dv/dr \sim v/r \sim 200~\text{km/s} / 10~\text{kpc} \sim 6 \times
    10^{-16}$~s$^{-1}$ is approximately 100 times larger than the Hubble
    flow gradient $\nabla \cdot \mathbf{v} = 3H \sim 7 \times
    10^{-18}$~s$^{-1}$.\ Viscous dissipation is proportional to the square
    of the velocity gradient, so the heating rate near mass can dominate
    over the heating rate in voids by a factor of $\sim$$10^4$.

  \item \textbf{Superfluid helium-4 is stiffer when cold.}\ The first-sound
    speed in He-4 is higher at lower temperatures: cold superfluid has a
    larger effective bulk modulus.\ If the Aether shares this property, cold
    voids would have higher $\Ka$ and higher wave speed $c$ than warm
    matter-containing regions.
\end{enumerate}

The resulting model---warm near mass, cold in voids, with
temperature-dependent $\Ka$---produces both dark matter and dark energy
from a single mechanism, governed by one parameter.\ The remainder of this
paper develops the quantitative framework and checks its consistency with
observations.


% ============================================================
\section{The Hot Matter Model}
\label{sec:hot-matter}
% ============================================================

\subsection{Temperature Reversal}

In the standard Paper~6 model, permanent binding continuously removes thermal
energy from the Aether near matter, keeping mass-containing regions cold.\
The alternative proposed here is that this thermal deficit has been more than
compensated by viscous heating from gravitational flows.

The normal (viscous) component of the two-fluid Aether dissipates kinetic
energy from any flow with velocity gradients.\ The dissipation rate per unit
volume is:
%
\begin{equation}
  \dot\varepsilon \;=\; \zeta_n\,(\nabla \cdot \mathbf{v})^2
    \;+\; \mu_n\left(\frac{\partial v_i}{\partial x_j}
    + \frac{\partial v_j}{\partial x_i}
    - \tfrac{2}{3}\delta_{ij}\nabla \cdot \mathbf{v}\right)^2
  \label{eq:dissipation-general}
\end{equation}
%
Near galaxies, gravitational inflow produces both shear and compressional
velocity gradients of order $v/r \sim 6 \times 10^{-16}$~s$^{-1}$ (for a
Milky-Way-sized galaxy at 10~kpc).\ In cosmic voids, the Hubble flow is
pure dilation with $\nabla \cdot \mathbf{v} = 3H \sim 7 \times 10^{-18}$~s$^{-1}$
and zero shear.\ The ratio of dissipation rates scales as the square of the
velocity gradients:
%
\begin{equation}
  \frac{\dot\varepsilon_{\mathrm{mass}}}{\dot\varepsilon_{\mathrm{void}}}
    \;\sim\; \left(\frac{6 \times 10^{-16}}{7 \times 10^{-18}}\right)^2
    \;\sim\; 10^4
  \label{eq:heating-ratio}
\end{equation}
%
Even after accounting for the lower normal fraction near mass (which reduces
$\zeta_n$ and $\mu_n$), the gravitational flow gradients are so much larger
that the near-mass heating rate can dominate.


\subsection{Self-Reinforcing Feedback}

The reversed temperature landscape is maintained by a positive feedback loop
on both sides:

\begin{enumerate}[leftmargin=2em]
  \item \textbf{Near mass (warm).}\ Higher temperature $\Rightarrow$ larger
    normal fraction $\rho_n/\rhoa$ $\Rightarrow$ larger bulk and shear
    viscosity $\Rightarrow$ more dissipation from gravitational inflow
    $\Rightarrow$ more heating $\Rightarrow$ stays warm.

  \item \textbf{In voids (cold).}\ Lower temperature $\Rightarrow$ smaller
    normal fraction $\Rightarrow$ lower viscosity $\Rightarrow$ negligible
    dissipation (only the weak Hubble-flow dilation contributes, and even that
    is suppressed by the low $\rho_n$) $\Rightarrow$ expansion cools the void
    further $\Rightarrow$ stays cold.
\end{enumerate}

This is a stable configuration: neither side has a mechanism to reverse its
thermal state.\ The warm region cannot cool (heating exceeds any losses),
and the cold region cannot warm (no significant heat source).\ By contrast,
the standard Paper~6 model must explain how voids remain hot despite
adiabatic cooling from expansion---a question addressed by Paper~6(a)'s
roton freeze-out mechanism \cite{otto-cooling}.


\subsection{Temperature-Dependent Bulk Modulus}
\label{sec:KA-T}

In superfluid helium-4, the first-sound speed $c_1 = \sqrt{K/\rho}$ is
higher at lower temperatures.\ At $T = 0$, the superfluid is maximally
stiff; as temperature increases toward the lambda point ($T_\lambda =
2.17$~K), thermal excitations soften the medium and $c_1$ decreases.\ The
fractional sensitivity at low temperature ($T < 1$~K) is approximately
$(1/c_1^2)|dc_1^2/dT| \sim 10^{-4}$~K$^{-1}$.

If the Aether exhibits the same behavior, the bulk modulus depends on
temperature:
%
\begin{equation}
  \Ka(T) \;=\; K_0\bigl[1 - \etaK\,(T - T_{\mathrm{ref}})\bigr]
  \label{eq:KA-T}
\end{equation}
%
where $K_0 \sim 10^{28}$~J/m$^3$ is the bulk modulus at the reference
temperature, and
%
\begin{equation}
  \etaK \;\equiv\; \frac{1}{\Ka}\left|\frac{d\Ka}{dT}\right|
  \label{eq:eta-def}
\end{equation}
%
is the fractional temperature sensitivity (units: K$^{-1}$).\ The wave speed
becomes position-dependent through the temperature field:
%
\begin{equation}
  c(r)^2 \;=\; \frac{\Ka(T(r))}{\rhoa}
  \label{eq:c-of-r}
\end{equation}
%
In the reversed model: near mass (warm, higher $T$) $\Rightarrow$ lower $\Ka$
$\Rightarrow$ lower $c$.\ In voids (cold, lower $T$) $\Rightarrow$ higher
$\Ka$ $\Rightarrow$ higher $c$.\ The wave speed \emph{increases outward} from
galaxies---precisely the gradient needed for the dark matter effect in the
acoustic metric.

The equation of state becomes $P = \Ka(T)\ln(\rho/\rho_0)$, coupling
temperature to mechanical pressure.\ This is a departure from Paper~1's
strict decoupling of the three properties, but the coupling is extremely
weak: $\Delta\Ka/\Ka \sim 10^{-6}$ on galactic scales and $\sim$$10^{-4}$
on cosmic scales.\ Lorentz invariance, the Fresnel drag coefficient, and all
results derived from constant $\Ka$ hold to better than one part in $10^4$.


% ============================================================
\section{Dark Matter from \texorpdfstring{$\Ka(T)$}{K\_A(T)}}
\label{sec:dark-matter}
% ============================================================

\subsection{Acoustic Metric Mechanism}

The orbital equation in the acoustic metric (Paper~6, Eq.~6) generalizes
to position-dependent wave speed:
%
\begin{equation}
  v_{\mathrm{orbit}}^2 \;=\; \frac{GM}{r}
    \;+\; \frac{r}{2}\,\frac{dc^2}{dr}
  \label{eq:vorbit}
\end{equation}
%
The dark matter contribution is the second term.\ Using the chain rule:
%
\begin{equation}
  \frac{dc^2}{dr} \;=\; \frac{1}{\rhoa}\,\frac{d\Ka}{dr}
    \;=\; \frac{K_0\,\etaK}{\rhoa}\left(-\frac{dT}{dr}\right)
  \label{eq:dc2dr}
\end{equation}
%
Since temperature \emph{decreases} outward ($dT/dr < 0$), the wave speed
\emph{increases} outward ($dc^2/dr > 0$), producing additional centripetal
acceleration---the dark matter effect with the correct sign.


\subsection{Flat Rotation Curves}

For a flat rotation curve ($v_{\mathrm{orbit}} = v_{\mathrm{flat}} =
\mathrm{const}$) at large $r$ where the dark matter term dominates:
%
\begin{equation}
  v_{\mathrm{flat}}^2 \;\approx\; \frac{r}{2}\,\frac{dc^2}{dr}
    \;=\; \frac{K_0\,\etaK}{2\,\rhoa}\,r\left|
    \frac{dT}{dr}\right|
  \label{eq:flat-rotation}
\end{equation}
%
For this to be independent of $r$, we require $r\,|dT/dr| = \mathrm{const}$,
which gives:
%
\begin{equation}
  T(r) \;=\; T_{\mathrm{center}} - \frac{2\,\rhoa\,v_{\mathrm{flat}}^2}
    {K_0\,\etaK}\,\ln\!\left(\frac{r}{r_{\mathrm{core}}}\right)
  \label{eq:T-profile}
\end{equation}
%
A logarithmic temperature profile is the natural steady-state solution for
viscous heating from a central gravitational flow source in a medium with
constant thermal conductivity.


\subsection{Quantitative Constraints}

The total $\Ka$ variation across a halo (from $r_{\mathrm{core}} = 1$~kpc
to $r_{\mathrm{halo}} = 200$~kpc, with $v_{\mathrm{flat}} = 220$~km/s) is:
%
\begin{equation}
  \frac{\Delta\Ka}{\Ka} \;=\; \frac{2\,\rhoa\,v_{\mathrm{flat}}^2}{\Ka}
    \,\ln\!\left(\frac{r_{\mathrm{halo}}}{r_{\mathrm{core}}}\right)
  \;=\; \frac{2 \times (220~\text{km/s})^2}{c^2}\,\ln(200)
  \;\approx\; 5.7 \times 10^{-6}
  \label{eq:DKA-halo}
\end{equation}
%
This result is \emph{independent of $\etaK$}: the fractional wave speed
variation required for flat rotation curves depends only on
$v_{\mathrm{flat}}/c$ and the logarithmic extent of the halo.\ The parameter
$\etaK$ determines how this variation maps to a temperature difference:
%
\begin{equation}
  \Delta T_{\mathrm{halo}} \;=\; \frac{\Delta\Ka/\Ka}{\etaK}
  \label{eq:DT-halo}
\end{equation}

\begin{table}[h]
\centering
\begin{tabular}{@{}lcc@{}}
\toprule
$\etaK$ (K$^{-1}$) & $\Delta T_{\mathrm{halo}}$ (K) &
  Comparison with He-4 \\
\midrule
$10^{-5}$        & 0.57     & $10\times$ gentler than He-4 at low $T$ \\
$5 \times 10^{-5}$ & 0.11   & $2\times$ gentler \\
$1.07 \times 10^{-4}$ & 0.053 & \textbf{Same as He-4 at low $T$} \\
$5 \times 10^{-4}$ & 0.011 & $5\times$ steeper \\
\bottomrule
\end{tabular}
\caption{The $\etaK$--$\Delta T_{\mathrm{halo}}$ degeneracy.\ For helium-4
  below 1~K, $(1/c_1^2)|dc_1^2/dT| \sim 10^{-4}$~K$^{-1}$.\ The benchmark
  $\etaK = 1.07 \times 10^{-4}$~K$^{-1}$ matches this value exactly.}
\label{tab:eta-DT}
\end{table}


\subsection{Gravitational Lensing}

Photons propagate at $c(r) = \sqrt{\Ka(T(r))/\rhoa}$.\ The same $c(r)$
gradient that produces enhanced rotation curves also bends light through the
acoustic metric.\ The effective mass profile is the singular isothermal
sphere (SIS):
%
\begin{equation}
  M_{\mathrm{eff}}(r) \;=\; \frac{v_{\mathrm{flat}}^2\,r}{G}
  \label{eq:Meff}
\end{equation}
%
and the deflection angle is:
%
\begin{equation}
  \theta \;=\; \frac{4\pi\,\sigma^2}{c^2}
  \label{eq:deflection}
\end{equation}
%
where $\sigma = v_{\mathrm{flat}}/\sqrt{2}$.\ This is identical to the
Paper~6 result: rotation curves and gravitational lensing are automatically
consistent because both light and matter follow geodesics of the same
acoustic metric.


\subsection{Profile Shape and Cores}

The logarithmic temperature profile (Eq.~\ref{eq:T-profile}) produces flat
rotation at large $r$, consistent with isothermal-sphere dark matter halos.\
At $r < r_{\mathrm{core}}$, the temperature profile flattens: converging
gravitational flows produce canceling velocity gradients, reducing the net
shear and therefore the viscous heating rate.\ As $dT/dr \to 0$, the dark
matter contribution vanishes, producing \emph{naturally cored} profiles
without additional assumptions---the same qualitative behavior as the
first-sound mechanism of Paper~6.

\emph{Galaxy clusters.}\ For a massive cluster with velocity dispersion
$\sigma \sim 1000$~km/s, the required $\Delta\Ka/\Ka \approx 10^{-4}$,
roughly $17\times$ larger than for a Milky-Way-sized galaxy.\ This is
physically expected: clusters are more massive, with stronger gravitational
inflow and correspondingly larger viscous heating---producing a larger
temperature gradient and therefore a larger $\Ka$ variation.


\subsection{Galactic Temperature Profile}
\label{sec:gal-profile}

Using the benchmark $\etaK = 1.07 \times 10^{-4}$~K$^{-1}$ and
anchoring the halo edge to $T_{\mathrm{CMB}} = 2.725$~K
(Section~\ref{sec:cmb}):

\begin{table}[h]
\centering
\begin{tabular}{@{}lccc@{}}
\toprule
Location & $r$ & $T$ (K) & $\Delta c/c$ \\
\midrule
Galaxy center    & 0         & 2.778 & $-2.8 \times 10^{-6}$ \\
Core edge        & 1~kpc     & 2.778 & (flat core) \\
Solar radius     & 8~kpc     & 2.747 & $-1.2 \times 10^{-6}$ \\
Disk edge        & 15~kpc    & 2.742 & $-0.9 \times 10^{-6}$ \\
Mid-halo         & 50~kpc    & 2.733 & $-0.4 \times 10^{-6}$ \\
Halo edge        & 200~kpc   & 2.725 & 0 (reference) \\
\bottomrule
\end{tabular}
\caption{Temperature profile across a Milky-Way-sized galaxy in the hot matter
  model, with $v_{\mathrm{flat}} = 220$~km/s, $\etaK = 1.07 \times
  10^{-4}$~K$^{-1}$, and halo edge anchored to $T_{\mathrm{CMB}}$.\ The total
  variation is only 0.053~K across 200~kpc.\ $\Delta c/c$ is relative to the
  halo edge.}
\label{tab:gal-profile}
\end{table}

The entire galaxy sits within a $\sim$0.05~K thermal envelope around
$T_{\mathrm{CMB}}$.\ This extraordinarily gentle gradient is sufficient
for flat rotation curves because even a fractional wave speed variation
of $\Delta c/c \sim 10^{-6}$ produces the full dark matter effect through
the acoustic metric.


% ============================================================
\section{Dark Energy from the Cooling Cascade}
\label{sec:dark-energy}
% ============================================================

\subsection{\texorpdfstring{$\Ka$}{K\_A} Gradient on Cosmic Scales}

On cosmic scales, the $\Ka$ contrast between warm walls
($T_{\mathrm{wall}} \approx 2.725$~K) and cold void centers
($T_{\mathrm{void}} \approx 0.91$~K) produces an outward acceleration
through the acoustic metric:
%
\begin{equation}
  a_K \;=\; \frac{1}{2}\,\frac{dc^2}{dr}
    \;\approx\; \frac{1}{2\,\rhoa}\,\frac{\Delta\Ka}{R_{\mathrm{void}}}
  \label{eq:aK}
\end{equation}
%
The observed dark energy acceleration at the boundary of a typical void
($R_{\mathrm{void}} \sim 50$~Mpc):
%
\begin{equation}
  a_{\mathrm{DE}} \;=\; R_{\mathrm{void}}\,H_0^2\,\Omega_\Lambda
    \;\approx\; 5.7 \times 10^{-12}~\text{m/s}^2
  \label{eq:aDE}
\end{equation}
%
Setting $a_K = a_{\mathrm{DE}}$ gives:
%
\begin{equation}
  \frac{\Delta\Ka_{\mathrm{cosmic}}}{\Ka}
    \;=\; \frac{2\,\rhoa\,R_{\mathrm{void}}\,a_{\mathrm{DE}}}{\Ka}
    \;\approx\; 1.95 \times 10^{-4}
  \label{eq:DKA-cosmic}
\end{equation}
%
With $\etaK = 1.07 \times 10^{-4}$~K$^{-1}$:
%
\begin{equation}
  \Delta T_{\mathrm{cosmic}} \;=\;
    \frac{\Delta\Ka_{\mathrm{cosmic}}/\Ka}{\etaK}
    \;\approx\; 1.82~\text{K}
  \label{eq:DT-cosmic}
\end{equation}
%
A temperature difference of less than 2~K between walls and void centers
is sufficient to produce the observed dark energy.\ Both the dark matter
constraint ($\Delta\Ka/\Ka \sim 6 \times 10^{-6}$ on galactic scales) and
the dark energy constraint ($\Delta\Ka/\Ka \sim 2 \times 10^{-4}$ on cosmic
scales) are satisfied by the \emph{same} $\etaK$.


\subsection{The Cooling Cascade}

The $\Ka$ gradient is not static---it is maintained and strengthened by
a self-reinforcing feedback loop:

\begin{enumerate}[leftmargin=2em]
  \item \textbf{Cold void has high $\Ka$.}\ The wave speed $c$ is elevated,
    and the acoustic metric drives expansion of the void.

  \item \textbf{Expansion cools the void.}\ Adiabatic cooling:
    %
    \begin{equation}
      \frac{dT_{\mathrm{void}}}{dt} \;=\; -\beta\,H\,T_{\mathrm{void}}
      \label{eq:Tvoid-evolution}
    \end{equation}
    %
    where $\beta$ parameterizes the effective adiabatic index.

  \item \textbf{Cooling increases $\Ka$.}\ Lower temperature $\Rightarrow$
    higher $\Ka$ $\Rightarrow$ higher $c$ $\Rightarrow$ stronger expansion
    drive.

  \item \textbf{Return to step 1.}\ The cascade is self-reinforcing: each
    cycle of expansion produces further cooling, which increases $\Ka$,
    which drives more expansion.
\end{enumerate}

Meanwhile, matter-containing walls are maintained at approximately constant
temperature by the viscous heating feedback (Section~\ref{sec:hot-matter}).

This cascade is \emph{more direct} than the roton freeze-out mechanism of
Paper~6(a) \cite{otto-cooling}.\ In Paper~6(a), the feedback goes: cooling
$\to$ roton freeze-out $\to$ viscosity collapse $\to$ less braking $\to$
faster expansion (indirect, through viscosity).\ Here: cooling $\to$ $\Ka$
increase $\to$ $c$ increase $\to$ stronger expansion drive (direct, through
the wave speed).

The void temperature evolves as:
%
\begin{equation}
  T_{\mathrm{void}}(z) \;=\; T_{\mathrm{void},0}\,(1 + z)
  \label{eq:Tvoid-z}
\end{equation}
%
while $T_{\mathrm{wall}} \approx \mathrm{const}$.\ The $\Ka$ contrast at
redshift $z$ is:
%
\begin{equation}
  \frac{\Delta\Ka(z)}{\Ka} \;=\; \etaK\bigl[T_{\mathrm{wall}}
    - T_{\mathrm{void},0}\,(1 + z)\bigr]
  \label{eq:DKA-z}
\end{equation}


\subsection{Cascade Evolution}

The temperature contrast \emph{reverses} at:
%
\begin{equation}
  z_{\mathrm{rev}} \;=\; \frac{T_{\mathrm{wall}}}{T_{\mathrm{void},0}} - 1
  \label{eq:z-reversal}
\end{equation}
%
For $T_{\mathrm{wall}} = 2.725$~K and $T_{\mathrm{void},0} = 0.91$~K:
$z_{\mathrm{rev}} \approx 2.0$.\ At $z > 2$, voids were \emph{warmer} than
walls, and no dark energy operated.

\begin{table}[h]
\centering
\begin{tabular}{@{}lcccc@{}}
\toprule
$z$ & $T_{\mathrm{void}}$ (K) & $\Delta T$ (K) &
  $\Delta\Ka/\Ka$ & DE fraction \\
\midrule
0   & 0.91 & 1.82 & $1.95 \times 10^{-4}$ & 100\% \\
0.5 & 1.37 & 1.36 & $1.45 \times 10^{-4}$ & 74\% \\
1.0 & 1.82 & 0.91 & $0.97 \times 10^{-4}$ & 50\% \\
2.0 & 2.73 & 0.00 & 0 & 0\% \\
5.0 & 5.46 & $-2.74$ & reversed & 0\% \\
\bottomrule
\end{tabular}
\caption{Evolution of the $\Ka$ contrast.\ At $z > 2$, voids are warmer than
  walls and no dark energy operates.\ The contrast builds gradually as voids
  cool, reaching full strength at $z = 0$.}
\label{tab:cascade}
\end{table}

This provides a natural explanation for \emph{why dark energy is a
late-universe phenomenon}: the $\Ka$ contrast had to build up through
billions of years of void cooling before it became cosmologically significant.


% ============================================================
\section{The \texorpdfstring{$w(z)$}{w(z)} Equation of State}
\label{sec:wz}
% ============================================================

\subsection{Phantom--Quintessence Crossing}

The effective dark energy density is proportional to the $\Ka$ gradient,
which evolves through two competing effects:

\begin{itemize}[leftmargin=2em]
  \item \textbf{Cascade growth} (strengthens dark energy): as voids cool,
    $\Delta\Ka$ increases, adding to the expansion drive.
  \item \textbf{Expansion dilution} (weakens dark energy): as voids grow
    larger, the gradient $\Delta\Ka/R_{\mathrm{void}}$ decreases for
    fixed $\Delta\Ka$.
\end{itemize}

The equation of state parameter:
%
\begin{equation}
  w(z) \;=\; -1 \;-\; \frac{1}{3H}\,\frac{d\ln\rho_{\mathrm{DE}}}{dt}
  \label{eq:w-def}
\end{equation}
%
The cascade growth rate is proportional to $T_{\mathrm{void}}$ (faster
cooling when voids are warmer), while the expansion dilution rate is
proportional to $H$.\ The crossing ($w = -1$) occurs when these rates
balance:
%
\begin{equation}
  T_{\mathrm{void,cross}} \;=\;
    \frac{T_{\mathrm{wall}}}{1 + \beta}
  \label{eq:T-cross}
\end{equation}
%
For the DESI crossing at $z_{\mathrm{cross}} \sim 0.5$:
%
\begin{equation}
  T_{\mathrm{void},0}\,(1 + z_{\mathrm{cross}})
    \;=\; \frac{T_{\mathrm{wall}}}{1 + \beta}
  \label{eq:z-cross}
\end{equation}
%
With $T_{\mathrm{wall}} = 2.725$~K and $z_{\mathrm{cross}} = 0.5$:
$T_{\mathrm{void},0} = T_{\mathrm{wall}}/[1.5\,(1 + \beta)]$.

\begin{itemize}[leftmargin=2em]
  \item At $z > z_{\mathrm{cross}}$: cascade growth dominates
    $\Rightarrow$ $\rho_{\mathrm{DE}}$ increasing $\Rightarrow$ $w < -1$
    \textbf{(phantom)}.
  \item At $z < z_{\mathrm{cross}}$: expansion dilution dominates
    $\Rightarrow$ $\rho_{\mathrm{DE}}$ decreasing $\Rightarrow$ $w > -1$
    \textbf{(quintessence)}.
\end{itemize}

This is the phantom-to-quintessence crossing observed by DESI at
2.8--4.2$\sigma$ significance \cite{desi2025}.


\subsection{Constraining Parameters from DESI}

The crossing condition (Eq.~\ref{eq:z-cross}) combined with the dark energy
magnitude (Eq.~\ref{eq:DT-cosmic}) and the dark matter constraint
(Eq.~\ref{eq:DKA-halo}) yields a self-consistent parameter set:

\begin{table}[h]
\centering
\begin{tabular}{@{}lcc@{}}
\toprule
Parameter & $\beta = 1$ & $\beta = 3$ \\
\midrule
$T_{\mathrm{wall}}$        & 2.725~K & 2.725~K \\
$T_{\mathrm{void},0}$      & 0.91~K  & 0.45~K  \\
$\Delta T_{\mathrm{cosmic}}$ & 1.82~K & 2.27~K \\
$\etaK$                     & $1.07 \times 10^{-4}$ & $8.6 \times 10^{-5}$ \\
$\Delta T_{\mathrm{halo}}$  & 0.053~K & 0.066~K \\
$z_{\mathrm{rev}}$          & 2.0     & 5.1     \\
\bottomrule
\end{tabular}
\caption{Self-consistent parameter sets for $\beta = 1$ and $\beta = 3$,
  with $T_{\mathrm{wall}} = 2.725$~K and DESI crossing at
  $z_{\mathrm{cross}} = 0.5$.}
\label{tab:parameters}
\end{table}

Both solutions give $\etaK \sim 10^{-4}$~K$^{-1}$, matching the measured
temperature sensitivity of helium-4's first-sound speed at low temperature.


\subsection{Comparison with Paper 6(a)}

Paper~6(a) and this paper predict the same qualitative $w(z)$
trajectory---phantom at high $z$, quintessence today, with a crossing near
$z \sim 0.5$--$1$---but through different physical mechanisms:

\begin{table}[h]
\centering
\begin{tabular}{@{}lll@{}}
\toprule
Feature & Paper~6(a) & This paper \\
\midrule
Temperature near mass & Cold & Warm \\
Temperature in voids  & Hot  & Cold \\
Dark matter mechanism & $c_1$ gradient (two-fluid) & $c$ gradient ($\Ka(T)$) \\
Dark energy mechanism & Viscosity collapse & $\Ka$ cascade \\
Feedback pathway      & Indirect (via $\zeta_n$) & Direct (via $\Ka \to c$) \\
Key parameter         & $\Delta/k_B \approx 6$~K & $\etaK \sim 10^{-4}$~K$^{-1}$ \\
He-4 analog           & Roton gap & Sound speed sensitivity \\
\bottomrule
\end{tabular}
\caption{Comparison of the two alternative dark sector mechanisms.}
\label{tab:comparison}
\end{table}


% ============================================================
\section{CMB and Local Temperature}
\label{sec:cmb}
% ============================================================

\subsection{\texorpdfstring{$T_{\mathrm{CMB}}$}{T\_CMB} at Our Position}

Paper~6 identifies the CMB as the thermal excitation spectrum of the Aether:
``$T_{\mathrm{CMB}} = 2.725$~K is the Aether's temperature''
\cite{otto-cosmology}.\ In the hot matter model, the Aether temperature
at the Sun's galactic position (8~kpc from center) is:
%
\begin{equation}
  T_\odot \;\approx\; 2.747~\text{K}
  \label{eq:Tsun}
\end{equation}
%
(Table~\ref{tab:gal-profile})---only 22~mK above $T_{\mathrm{CMB}}$, a
0.8\% offset.

This near-equality is not a coincidence.\ Both $T_{\mathrm{CMB}}$ and
$T_{\mathrm{wall}}$ are governed by the same cosmic expansion history:
the CMB temperature tracks the relic photon field redshifted from last
scattering ($T_{\mathrm{CMB}} = T_{\mathrm{LSS}}/(1 + z_{\mathrm{LSS}})
\approx 2.725$~K), while the wall temperature is set by the balance of
viscous heating and expansion cooling---both of which depend on $H(t)$ and
the age of the universe.\ The two converge to similar values at the current
epoch.

The hot matter model has a structural advantage here.\ In the standard
Paper~6 model, $T_{\mathrm{wall}} \approx 2.7$~K but
$T_{\mathrm{void}} \approx 7$~K, giving a volume-weighted average of
$0.3 \times 2.7 + 0.7 \times 7 \approx 5.7$~K---not $T_{\mathrm{CMB}}$.\
In the hot matter model, the locally measured temperature (at our position
near mass) \emph{is} $T_{\mathrm{CMB}}$ to within 1\%.


\subsection{CMB Isotropy}

The galactic temperature gradient (0.053~K across 200~kpc) would produce
$\delta T/T \sim 10^{-2}$ on galactic angular scales if the CMB reflected
the local Aether temperature.\ However, CMB photons from the last scattering
surface propagate through the local Aether without thermalizing: direct
photon--quasiparticle scattering is suppressed by
$P_{\mathrm{th}}/\Ka \sim 10^{-38}$ \cite{otto-cosmology}.\ The observed
anisotropy ($\delta T/T \sim 10^{-5}$) reflects the primordial temperature
fluctuations at $z \sim 1100$, not the current local temperature structure.\
The $\Ka$ gradient contributes to the integrated Sachs--Wolfe effect
(photons gaining or losing energy as they traverse evolving $c(r)$ regions),
at a level consistent with observations.


% ============================================================
\section{Observational Consistency}
\label{sec:observations}
% ============================================================

\subsection{Rotation Curves}

\textbf{Observation:} Flat rotation curves with $v_{\mathrm{flat}} \sim
200$~km/s for Milky-Way-sized galaxies.

\textbf{Model:} $\Delta\Ka/\Ka \approx 5.7 \times 10^{-6}$ from the
logarithmic temperature profile (Eq.~\ref{eq:T-profile}).

\textbf{Status:} \textbf{Consistent by construction.}


\subsection{Gravitational Lensing}

\textbf{Observation:} Lensing mass agrees with dynamical mass from rotation
curves.

\textbf{Model:} Both light and matter follow geodesics of the same acoustic
metric with position-dependent $c(r)$.\ No separate mechanism needed.

\textbf{Status:} \textbf{Automatically consistent.}


\subsection{Cored Dark Matter Profiles}

\textbf{Observation:} Dwarf and low-surface-brightness galaxies show cored
(not cuspy) dark matter profiles.

\textbf{Model:} The temperature profile flattens at $r < r_{\mathrm{core}}$
(converging flows reduce velocity gradients), so $dT/dr \to 0$ and the dark
matter contribution vanishes smoothly.

\textbf{Status:} \textbf{Naturally cored.}


\subsection{Bullet Cluster}

\textbf{Observation:} Gravitational lensing follows the galaxies, not the
shocked gas, in the 1E~0657-56 cluster collision.

\textbf{Model:} Each galaxy's temperature profile (and therefore $\Ka(r)$
gradient) is maintained by its own gravitational inflow.\ The shocked gas
between the sub-clusters heats the local Aether, \emph{reducing} $\Ka$ and
weakening the $c$ gradient in the collision region.\ The surviving $\Ka$
gradients follow the galaxies, not the gas---producing the observed
lensing--gas separation.

\textbf{Status:} \textbf{Consistent.}


\subsection{Hubble Tension}

\textbf{Observation:} Local measurements give $H_0 \approx 73$~km\,s$^{-1}$\,Mpc$^{-1}$;
CMB-derived values give $\approx$67~km\,s$^{-1}$\,Mpc$^{-1}$.

\textbf{Model:} The KBC void (local underdensity, $\delta \approx 0.25$)
is colder than average $\Rightarrow$ higher $\Ka$ $\Rightarrow$ higher $c$.\
Photons from distant sources traversing this region travel at higher-than-average
speed, biasing local distance measurements and producing a higher apparent $H_0$.\
The qualitative mechanism parallels Paper~6's thermal pressure excess in the void
\cite{otto-cosmology}, but operates through $\Ka$ rather than $P_{\mathrm{th}}$.

\textbf{Status:} \textbf{Qualitatively consistent; quantitative calculation
needed.}


\subsection{Galaxy Clusters}

\textbf{Observation:} Clusters require $\Delta c^2/c^2 \sim 10^{-4}$,
roughly $17\times$ larger than galaxies.

\textbf{Model:} Clusters have more massive gravitational inflows, stronger
velocity gradients, and therefore larger viscous heating---producing a larger
temperature contrast and $\Ka$ variation.\ The scaling
$\Delta\Ka/\Ka \propto v_{\mathrm{flat}}^2$ naturally produces the observed
galaxy-to-cluster ratio.

\textbf{Status:} \textbf{Consistent.}


\subsection{DESI Dark Energy Equation of State}

\textbf{Observation:} Phantom-to-quintessence crossing at $z \gtrsim 0.5$
\cite{desi2025}.

\textbf{Model:} The cooling cascade produces $w < -1$ at high $z$ (cascade
growth dominates) transitioning to $w > -1$ at low $z$ (expansion dilution
dominates), with crossing at $z \sim 0.5$ for $\beta = 1$
(Table~\ref{tab:parameters}).

\textbf{Status:} \textbf{Consistent.}


% ============================================================
\section{Predictions}
\label{sec:predictions}
% ============================================================

\begin{enumerate}[leftmargin=2em]
  \item \textbf{$\etaK \sim 10^{-4}$~K$^{-1}$.}\ The temperature
    sensitivity of the Aether's bulk modulus must match the measured
    sensitivity of helium-4 at low temperature.\ If the Aether's excitation
    spectrum can be derived from first principles, this is a specific,
    testable prediction.

  \item \textbf{$w(z)$ follows a cascade curve, not a polynomial.}\ The
    dark energy evolution is determined by the void cooling history
    (Eq.~\ref{eq:DKA-z}), not the smooth CPL parameterization.\ Future
    DESI data releases and Euclid could distinguish the two functional
    forms.

  \item \textbf{Dark energy is absent at $z > 2$.}\ The temperature
    contrast reverses (Table~\ref{tab:cascade}): at $z \gtrsim 2$, voids
    were warmer than walls, so no $\Ka$-driven dark energy operated.\ This
    contrasts with $\Lambda$CDM, where dark energy is present at all epochs.

  \item \textbf{$T_{\mathrm{CMB}} \approx T_{\mathrm{wall}}$.}\ The CMB
    temperature directly measures the Aether temperature in mass-containing
    regions (where all observers reside).\ Any departure from
    $T_{\mathrm{wall}} \approx T_{\mathrm{CMB}}$ would falsify this model.

  \item \textbf{Void regions are cold and expand faster.}\ Voids with
    lower Aether temperatures should show stronger local expansion rates and
    enhanced ISW imprints.\ This is testable with upcoming void catalogs
    from DESI and Euclid.

  \item \textbf{Single parameter determines both dark phenomena.}\ The same
    $\etaK$ that produces dark matter on galactic scales produces dark energy
    on cosmic scales.\ Independent measurements of both should converge on
    the same value---a powerful consistency test.
\end{enumerate}


% ============================================================
\section{Open Questions}
\label{sec:open}
% ============================================================

\begin{enumerate}[leftmargin=2em]
  \item \textbf{Derivation of $\etaK$ from Aether microscopic parameters.}\
    Can $\etaK \sim 10^{-4}$~K$^{-1}$ be derived from $m_A$, $\rhoa$, and
    $\Ka$?\ In helium-4, the temperature dependence of the sound speed
    arises from the phonon-roton excitation spectrum; the analogous
    calculation for the Aether requires knowledge of the full excitation
    spectrum.

  \item \textbf{Three-dimensional temperature profile.}\ The logarithmic
    $T(r)$ profile (Eq.~\ref{eq:T-profile}) produces flat rotation, but the
    full three-dimensional profile---including the transition from halo to
    void---needs to be derived from the viscous heating equation with
    appropriate boundary conditions.

  \item \textbf{Tully--Fisher relation.}\ The model predicts
    $v_{\mathrm{flat}}^2 \propto (\text{heating rate}) \times
    \ln(r_{\mathrm{halo}}/r_{\mathrm{core}})$.\ Whether this reproduces the
    observed $M \propto v^4$ scaling depends on how the viscous heating rate
    scales with baryonic mass.\ For $\dot\varepsilon \propto M^{0.7}$
    (sublinear, from the star-formation main sequence), the predicted slope
    approaches the observed value.

  \item \textbf{CMB power spectrum.}\ How does the $\Ka(T)$ mechanism affect
    baryon-photon acoustic oscillations at $z \sim 1100$?\ At the CMB epoch,
    the temperature was nearly uniform ($\delta T/T \sim 10^{-5}$), so
    $\delta\Ka/\Ka \sim \etaK \times \delta T \sim 10^{-4} \times 10^{-5}
    \times 3000 \sim 3 \times 10^{-6}$---small but potentially detectable in
    the damping tail.

  \item \textbf{Reconciliation with Paper~6.}\ The standard model (cold near
    mass) and the hot matter model (warm near mass) produce qualitatively
    similar observational signatures through different mechanisms.\ Are they
    limiting cases of a more general framework?\ Could the truth lie between
    the two extremes?

  \item \textbf{The $\beta$ parameter.}\ The adiabatic cooling rate of void
    Aether determines the cascade evolution and the $w(z)$ crossing redshift.\
    In a standard fluid, $\beta = 3$; for the nearly incompressible Aether,
    $\beta$ depends on how thermal energy couples to the expansion and
    requires a first-principles calculation.

  \item \textbf{Specific heat.}\ The viscous heating rate constrains the
    product $\rhoa\,c_v\,T_{\mathrm{wall}}$.\ For
    $T_{\mathrm{wall}} \approx 2.7$~K, the required specific heat is
    $c_v \sim 5 \times 10^{-20}$~J/(kg\,K)---consistent with a superfluid
    far below $T_c$ (Debye $c_v \propto (T/T_c)^3$), but the consistency
    with the same excitation spectrum that determines $\etaK$ needs to
    be verified.

  \item \textbf{Hubble tension magnitude.}\ The qualitative mechanism
    (cold KBC void $\Rightarrow$ higher $\Ka$ $\Rightarrow$ higher local
    $c$) is clear, but the quantitative requirement ($\sim$8\% discrepancy
    in $H_0$) needs detailed modeling of the KBC void temperature profile
    and its effect on distance-ladder measurements.
\end{enumerate}


% ============================================================
% REFERENCES
% ============================================================
\begin{thebibliography}{99}

\bibitem{desi2025}
DESI Collaboration,
``DESI DR2 Results: Measurement of the Baryon Acoustic Oscillation
scale and Hubble parameter across 11 billion years,''
\textit{arXiv:2503.14738} (2025).

\bibitem{berezhiani2015}
L.~Berezhiani and J.~Khoury,
``Theory of dark matter superfluidity,''
\textit{Phys.\ Rev.\ D}, \textbf{92}, 103510 (2015).

\bibitem{planck2020}
Planck Collaboration,
``Planck 2018 results.\ VI.\ Cosmological parameters,''
\textit{Astron.\ Astrophys.}, \textbf{641}, A6 (2020).

\bibitem{riess2022}
A.~G.~Riess \textit{et al.},
``A comprehensive measurement of the local value of the Hubble constant
with 1 km/s/Mpc uncertainty from the Hubble Space Telescope and the SH0ES
Team,''
\textit{Astrophys.\ J.\ Lett.}, \textbf{934}, L7 (2022).

\bibitem{otto-overview}
E.~Otto,
``Unifying Theory of Dynamic Aether Mechanics: General Overview,''
this series, Paper~1.

\bibitem{otto-gravity}
E.~Otto,
``Unifying Theory of Dynamic Aether Mechanics: Gravity,''
this series, Paper~3.

\bibitem{otto-cosmology}
E.~Otto,
``Unifying Theory of Dynamic Aether Mechanics: Cosmology,''
this series, Paper~6.

\bibitem{otto-cooling}
E.~Otto,
``Superfluid Cooling Feedback as the Origin of Dark Energy and Dark Matter,''
this series, Paper~6(a).

\end{thebibliography}

\end{document}
